\documentclass[UTF8,a4paper,12pt]{ctexart}

\usepackage{amsmath,amssymb,amsthm,mathtools}
\usepackage{geometry}
\usepackage{booktabs}
\usepackage{array}
\usepackage{enumitem}
\usepackage{hyperref}
\usepackage{titlesec}
\usepackage{fancyhdr}

\geometry{margin=1in}
\hypersetup{colorlinks=true,linkcolor=black,urlcolor=blue,citecolor=black}
\setlist[itemize]{leftmargin=2em}
\setlist[enumerate]{leftmargin=2.2em}

\title{不定积分的“双元法”整理}
\author{}
\date{}

\begin{document}

\maketitle
\tableofcontents
\newpage

\section{双元法核心思想}

双元法的本质可以概括为：

\begin{quote}
  构造两个变量 $x,y$，让它们满足某种平方关系，从而把积分化为便于处理的“凑微分”形式。
\end{quote}

常见关系为
\[
  x^2+y^2=\text{常数},\qquad x^2-y^2=\text{常数}.
\]

对应的微分关系分别为
\[
  xdx+ydy=0,\qquad xdx-ydy=0.
\]

双元法的作用主要体现在以下几点：
\begin{itemize}
  \item 将复杂根式、三角形式或对称结构转化为标准积分；
  \item 常将原式化为 $\int \frac{dx}{y}$ 或 $\int y\,dx$ 的形式；
  \item 本质上仍然属于“换元 + 凑微分”的方法。
\end{itemize}

\section{双元法三大基本公式}

设 $x,y$ 满足以下两类关系之一：
\[
  xdx+ydy=0 \quad \text{（对应平方和）}
\]
或
\[
  xdx-ydy=0 \quad \text{（对应平方差）}.
\]

\subsection{公式一：$\int \frac{dx}{y}$}

\[
  \int \frac{dx}{y}=
  \begin{cases}
    \arctan\dfrac{x}{y}+C, & xdx+ydy=0,\\[1em]
    \ln|x+y|+C, & xdx-ydy=0.
  \end{cases}
\]

这是双元法中最核心的入口公式。

\subsection{公式二：$\int \frac{dx}{y^3}$}

\[
  \int \frac{dx}{y^3}=\frac{1}{x^2\pm y^2}\cdot \frac{x}{y} + C.
\]

该公式常用于处理更高次的根式积分。

\subsection{公式三：$\int y\,dx$}

\[
  \int y\,dx=\frac{1}{2}xy+\frac{y^2\pm x^2}{2}\int \frac{dx}{y}.
\]

这个公式的本质可以看作分部积分与双元关系的结合。

\section{双元法的标准计算步骤}

解题时通常按如下步骤进行。

\subsection{Step 1：观察结构并构造双元}

常见识别方式如下：

\begin{center}
  \begin{tabular}{cc}
    \toprule
    类型 & 构造 \\
    \midrule
    $\sqrt{a^2-x^2}$ & $x^2+y^2=a^2$ \\
    $\sqrt{x^2+a^2}$ & $x^2-y^2=-a^2$ \\
    $\sqrt{x^2-a^2}$ & $x^2-y^2=a^2$ \\
    \bottomrule
  \end{tabular}
\end{center}

核心思想是：让根号中的表达式自然成为 $y$，从而消去根式。

\subsection{Step 2：写出微分关系}

例如由
\[
  x^2+y^2=a^2
\]
可得
\[
  xdx+ydy=0.
\]

这是整个方法中的关键桥梁。

\subsection{Step 3：把原积分改写为标准形式}

尽量把原式改写为下列两种常见类型之一：
\[
  \int \frac{dx}{y},\qquad \int y\,dx.
\]

这一步体现了“双元法 = 凑微分”的核心。

\subsection{Step 4：套用基本公式}

根据前述三大公式直接计算。

\subsection{Step 5：回代}

最后把辅助变量 $y$ 用原变量表示出来，得到最终结果。

\section{经典例题}

考虑积分
\[
  I=\int \frac{dx}{\sqrt{a^2-x^2}}.
\]

\subsection*{第一步：构造双元}
设
\[
  y=\sqrt{a^2-x^2},
\]
则有
\[
  x^2+y^2=a^2.
\]

\subsection*{第二步：写出微分关系}
由上式得
\[
  xdx+ydy=0.
\]

\subsection*{第三步：改写原积分}
于是
\[
  I=\int \frac{dx}{y}.
\]

\subsection*{第四步：套用公式一}
因为此时满足 $xdx+ydy=0$，故
\[
  I=\arctan\frac{x}{y}+C.
\]

\subsection*{第五步：回代}
代回 $y=\sqrt{a^2-x^2}$，得
\[
  I=\arctan\frac{x}{\sqrt{a^2-x^2}}+C.
\]

\section{常见适用题型}

双元法尤其适合以下类型的问题。

\subsection{根式积分}
\begin{itemize}
  \item $\sqrt{x^2+a^2}$
  \item $\sqrt{a^2-x^2}$
  \item $\sqrt{x^2-a^2}$
\end{itemize}

\subsection{对称结构}
\begin{itemize}
  \item $x+\sqrt{x^2+a}$
  \item $\dfrac{1}{x+\sqrt{x^2+1}}$
\end{itemize}

\subsection{三角或复合结构}
\begin{itemize}
  \item $\sqrt{\sin x-\sin^2 x}$
\end{itemize}

\section{考试中需要记住的本质}

一句话总结：

\begin{quote}
  双元法 = 构造关系 + 凑微分 + 套公式。
\end{quote}

\section{双元法与换元法的关系}

双元法本质上仍属于换元法，但它有自己的特点：
\begin{itemize}
  \item 结构更加对称；
  \item 操作更加模板化；
  \item 在部分根式问题中比三角换元更直接。
\end{itemize}

\section{公式一的推导：$\int \frac{dx}{y}$}

下面分别对两种情形进行推导。

\subsection{情形 A：$xdx+ydy=0$}

目标是证明
\[
  \int \frac{dx}{y}=\arctan\frac{x}{y}+C.
\]

先对 $\dfrac{x}{y}$ 求导：
\[
  d\!\left(\frac{x}{y}\right)=\frac{y\,dx-x\,dy}{y^2}.
\]

由双元关系
\[
  xdx+ydy=0
\]
得
\[
  x\,dx=-y\,dy.
\]

进而可写为
\[
  x\,dy=-\frac{x^2}{y}dx.
\]

代回导数公式：
\[
  d\!\left(\frac{x}{y}\right)=\frac{y\,dx-\left(-\frac{x^2}{y}dx\right)}{y^2}
  =\frac{y^2+x^2}{y^3}dx.
\]

若记
\[
  x^2+y^2=C,
\]
则有
\[
  d\!\left(\frac{x}{y}\right)=\frac{C}{y^3}dx.
\]

再利用
\[
  d(\arctan u)=\frac{du}{1+u^2},
\]
令 $u=\dfrac{x}{y}$，则
\[
  d\!\left(\arctan\frac{x}{y}\right)
  =\frac{d(x/y)}{1+(x/y)^2}.
\]

代入前式可得
\[
  d\!\left(\arctan\frac{x}{y}\right)
  =\frac{\frac{C}{y^3}dx}{\frac{x^2+y^2}{y^2}}
  =\frac{dx}{y}.
\]

因此
\[
  \int \frac{dx}{y}=\arctan\frac{x}{y}+C.
\]

\subsection{情形 B：$xdx-ydy=0$}

目标是证明
\[
  \int \frac{dx}{y}=\ln|x+y|+C.
\]

直接对 $x+y$ 求导：
\[
  d(x+y)=dx+dy.
\]

由
\[
  xdx-ydy=0
\]
得
\[
  ydy=xdx,
\]
所以
\[
  dy=\frac{x}{y}dx.
\]

代入得
\[
  d(x+y)=dx+\frac{x}{y}dx=\frac{x+y}{y}dx.
\]

于是
\[
  \frac{dx}{y}=\frac{d(x+y)}{x+y}.
\]

两边积分即得
\[
  \int \frac{dx}{y}=\int \frac{d(x+y)}{x+y}=\ln|x+y|+C.
\]

\section{公式三的推导：$\int y\,dx$}

目标证明
\[
  \int y\,dx=\frac{1}{2}xy+\frac{y^2\pm x^2}{2}\int \frac{dx}{y}.
\]

\subsection{分部积分思路}

先作分部积分，取
\[
  u=y,\qquad dv=dx.
\]
则
\[
  \int y\,dx=xy-\int x\,dy.
\]

再利用恒等变形
\[
  y\,dx=\frac{1}{2}(y\,dx+x\,dy)+\frac{1}{2}(y\,dx-x\,dy).
\]

注意到
\[
  y\,dx+x\,dy=d(xy).
\]

另一方面，由双元关系$xdx \pm ydy = 0$ 可以知道,
$ydy = \mp xdx$

\[
  y\,dx-x\,dy
\]
化为
\[
  \left(y^2\pm x^2\right)\frac{dx}{y}.
\]

于是积分两边得到
\[
  \int y\,dx
  =\frac{1}{2}\int d(xy)+\frac{1}{2}\int \left(y^2\pm x^2\right)\frac{dx}{y}.
\]

即
\[
  \int y\,dx
  =\frac{1}{2}xy+\frac{y^2\pm x^2}{2}\int \frac{dx}{y}.
\]

\section{最后总结}

双元法的全部关键都可以归结为以下三步：
\begin{enumerate}
  \item 建立双元关系：$xdx\pm ydy=0$；
  \item 构造合适微分：如 $d(x/y)$、$d(x+y)$、$d(xy)$；
  \item 把原积分化为标准形式：如 $\dfrac{dx}{y}$ 或 $y\,dx$。
\end{enumerate}

因此，双元法并不是“背公式”的技巧，而是通过变量关系把结构锁定，再制造出可积形式的方法。

\end{document}
